\scanpage{113}
% archive sheet number: 112
\textcircled{$2'$} Posons $U=X-Y$, qui est une variété affine non sing.
\[
 A^i_r(U)\simeq A^i_r(X)/\operatorname{Im} A^i_r(Y),\quad
 \text{de }=\text{ pour } A^i_a(U)\ldots
\]
\[
\begin{aligned}
 2i\leq n-2&\quad\text{i.e. }i\leq\frac n2-1&
 A_i^r(U)&=0\quad(\text{exc. $p$ ?})\deleted{\unsure{??}}\\
 2i=n-1&\quad\text{i.e. }i=\frac{n-1}2\deleted{\unsure{+1}}&
 A_i^a(U)&=0\quad(\text{exc $p$ ?})
\end{aligned}
\]

\textcircled{$3'$} Soit $E$ le cône projetant affine épointé de $X$, (variété quasi-affine non sing. de dim $2n+1$).

\[
\begin{aligned}
 2i\geq n+2&\quad\text{i.e. } i\geq\frac n2+1&
 A_r^i(E)&\ \deleted{\unsure{=0|\equiv}}\text{ fini},\quad A^i_r(E)^a=0\\
 &\quad n-i\leq\frac n2-1\\
 2i\geq n+1&\quad\text{i.e. } i\geq\frac{n+1}2&
 A^i_a(E)&\ \text{fini}\\
 &\quad\text{i.e. }\unsure{n|N}-i\leq\frac{n-1}2
\end{aligned}
\]

\uline{N.B.} On doit avoir
\[
\left\{\begin{aligned}
 A_a^i(E)&\simeq\deleted{\unsure{B_a^i(X)}}\unsure{N|M|P}_a^i(X)
 \quad\text{\inserted{mod groupes finis} si }2i\leq n\\
 \deleted{A_r^i(E)^a}&\deleted{\simeq\unsure{B_r^i|P^i_{\tau}}(X)\quad\text{mod gpes finis si }2i\leq n+1}\\
 A_a^i(E)&\simeq0\quad\text{mod gpes finis si }2i\geq n+1
\end{aligned}\right.
\qquad
\]
\[
\left\{\begin{aligned}
 A_r^i(E)^a&\simeq I^{i}_{\deleted{\unsure{B_r}}}(X)\quad\text{mod gpes finis si }2i\leq n+1\\
 A_r^i(E)^{\unsure{a|-}}&\simeq\unsure{0|B}\qquad\text{``}\qquad\text{si }2i\geq n+2
\end{aligned}\right.
\]
où
\[
\begin{aligned}
 \deleted{\unsure{B}}N^i_a(X)&\simeq\operatorname{Ker}\bigl(A^i_a(X)\xrightarrow{L^{n-2i}}A^{n-i}_{\unsure{a}}(X)\bigr)\\
 I^i_{\unsure{r|a}}(X)&=\operatorname{Ker}\bigl(A^i_r(X)^a\xrightarrow{L^{n-2i+1}}A^{n-i+1}_r(X)^a
\end{aligned}
\]
\margin{partie ``primitive'' de la ``\unsure{coh.} algébrique''}
\margin{partie ``primitive'' \deleted{de} la V.A. intermédiaire}
